\documentclass[12pt]{article}
\usepackage[utf8]{inputenc}
\usepackage{lmodern}
\usepackage{geometry}
\geometry{margin=2cm}

\author{Bakary keita}
\date{\today}
\title{Resolution numerique de l'équation de la diffusion}
\begin{document}
\maketitle 



Je traite le problème de la diffusion de la chaleur en 1D entre 0 et la
le schema de discretisation nous conduit aux expression obtenu dans [expre.ipynb]. 
En factorisant on trouve $Ti^{n+1}=GTi^n. Nous distinguons 2 cas: $G<=1$(satbilite) et $G>=1$
 (instabilité) avec $G=2r(1-4sin^2(kdx))$

Immédiate stable si $r>=1/2$.

Nous choisisson 4 points à l'interieur de mon domaine de définition c'est à dire pour 
$i=1;2;3;4$ une execution à la main conduit à  

\end{document}
